%% ----------------------------------------------------------------
%% uet_thesis.tex -- MAIN FILE (compile this file)
%% Unified FYP Design Document Template
%% ----------------------------------------------------------------

\documentclass[a4paper,12pt,oneside]{uet_thesis}
\graphicspath{{Figures/}}

% Packages required by the design-document template.
\usepackage[T1]{fontenc}
\usepackage[utf8]{inputenc}
\usepackage{lmodern}
\usepackage{microtype}
\usepackage[square,numbers,comma,sort&compress]{natbib}
\usepackage{url}
\usepackage{verbatim}
\usepackage{tabularx}
\usepackage{longtable}
\usepackage{array}
\usepackage{enumitem}
\usepackage{xcolor}
\usepackage{float}
\usepackage[compact]{titlesec}

% Match the compact heading treatment used by the supplied UET template.
\titlespacing{\section}{0pt}{*0}{*0}
\titlespacing{\subsection}{0pt}{*0}{*0}
\titlespacing{\subsubsection}{0pt}{*0}{*0}

\definecolor{GuidanceGray}{HTML}{4D4D4D}
\definecolor{PlaceholderBlue}{HTML}{1F5AA6}
\definecolor{ExampleGray}{HTML}{F2F4F7}

\hypersetup{urlcolor=blue,colorlinks=true,linkcolor=blue,citecolor=blue}
\pdfstringdefDisableCommands{\def\quad{ }}
\setlength{\parskip}{6pt}
\setlist[itemize]{topsep=0.25em,itemsep=0.15em,leftmargin=1.6em}
\setlist[enumerate]{topsep=0.25em,itemsep=0.15em,leftmargin=1.8em}
\renewcommand{\arraystretch}{1.2}

\newcommand{\placeholder}[1]{\textcolor{PlaceholderBlue}{\textit{<#1>}}}
\newcommand{\applicability}[1]{\texorpdfstring{\normalfont\small\textit{(#1)}}{(#1)}}

\newenvironment{guidance}
  {\begin{quote}\small\color{GuidanceGray}\textbf{Guidance.} }
  {\end{quote}}

\newenvironment{examplebox}
  {\begin{quote}\small\color{GuidanceGray}\textbf{Example (replace or remove).} }
  {\end{quote}}

\newcommand{\diagramplaceholder}[2]{%
  \begin{figure}[H]
    \centering
    \fcolorbox{GuidanceGray}{ExampleGray}{%
      \begin{minipage}[c][4.0cm][c]{0.86\textwidth}
        \centering
        \textbf{#1}\\[0.5em]
        \placeholder{Insert a clear, readable diagram here}\\[0.5em]
        \small #2
      \end{minipage}}
    \caption{#1}
  \end{figure}
}

\begin{document}
\frontmatter

% -----------------------------------------------------------------------------
% Title-page information: replace all sample values.
% -----------------------------------------------------------------------------
\title{\texorpdfstring{surplusSpace\\[1.5em]{\large Final Year Project Design Document}}{}}
\session{2023--2027 \quad Project Type: Development / Hybrid}
\advisor{Ms. Anam Iftikhar}
\authors{%
Hafsa Shahid \quad 2023-CS-708 \\
Fatima Irfan \quad 2023-CS-715 \\
Daniyal Jumshaid \quad 2023-CS-734 \\
Noor Fatima Malik \quad 2023-CS-739}
\addresses{\deptname \\ \univname}
\date{\today}
\subject{Final Year Project Design Document}
\keywords{software design, requirements, final year project}
\maketitle

\setstretch{1.3}
\fancyhead{}
\rhead{\thepage}
\lhead{}
\pagestyle{fancy}

\chapter*{Template Instructions}
\addcontentsline{toc}{chapter}{Template Instructions}
This template supports development, research, hybrid, and game projects. Complete
all common sections and retain conditional sections only when they apply to the
selected project. Keep the document concise and use cross-references instead of
repeating information.

\begin{itemize}
  \item Replace every \placeholder{placeholder} and remove all example text before submission.
  \item State the project type on the title page: Development, Research, Hybrid, or Game.
  \item Assign stable identifiers to objectives, scope items, requirements, risks, and design elements.
  \item Write one clear idea or testable requirement per numbered item.
  \item Use standard UML notation where UML diagrams are included.
  \item Cite every external idea, dataset, standard, figure, and factual claim.
  \item Keep detailed test cases and results in the separate software testing document; reference them through the traceability matrix.
\end{itemize}
\clearpage

\lhead{\emph{Design Document}}
\tableofcontents
\listoffigures
\listoftables

\mainmatter
\pagestyle{fancy}
\onehalfspacing

% This file is included by uet_thesis.tex.
% Replace blue placeholders and remove example text before submission.

\chapter{Project Overview}
\label{chap:overview}
% =============================================================================

\section{Executive Summary}
surplusSpace is an AI-gated marketplace for surplus construction materials. It
addresses the lack of a trustworthy channel through which sellers and buyers can
exchange surplus bricks, steel/rebar, tiles/marble, and wood. A seller creates a
lot with its quantity, description, price, and 3--5 representative photographs.
Before publication, the system classifies the material, detects selected visible
defects, derives a condition grade, checks description--image consistency, and
estimates a fair surplus price against a new-material reference price. A decision
engine automatically publishes listings that pass the checks and routes uncertain
or risky listings to an Admin, who has final authority to override the AI.
Customers and Contractors use a switchable account to search by location, price,
and category, negotiate through real-time chat, and reserve lots. An AI Listing
Assistant answers questions using only the listing's own data. The project is
limited to four material categories and evaluates classification, defect
detection, price estimation, and agreement between automated decisions and Admin
judgement.

\section{Background and Motivation}
Surplus bricks, steel/rebar, tiles/marble, and wood are often wasted because
sellers and buyers lack a trustworthy exchange channel. Existing classifieds
platforms do not automatically check listing condition, price fairness, or
description accuracy, while manual moderation is difficult to scale. The project
therefore combines a construction-material marketplace with an AI verification
step and a human Admin fallback. The scope is deliberately limited to four
material categories and documented defect classes so that the system can be
developed and evaluated within the 12-month project.

\section{Problem Statement}
Contractors and sellers listing surplus construction material have no automated
way to verify, before a listing reaches buyers, that its condition, price, and
description are consistent and honest. Existing marketplace platforms rely mainly
on manual or after-the-fact moderation, which does not scale and leaves buyers
exposed to misrepresented listings. A marketplace is therefore needed that
screens each listing before publication while retaining a human Admin as the
final authority.

\section{Research Questions and Hypotheses \applicability{Research, Hybrid; hypotheses if applicable}}
\begin{guidance}
State focused questions derived from the problem statement. Include hypotheses
only when the study tests predicted relationships or differences. Exploratory or
qualitative studies may not require hypotheses.
\end{guidance}

\begin{table}[H]
\centering
\caption{Research questions and hypotheses}
\label{tab:research-questions}
\begin{tabularx}{\textwidth}{@{}p{0.13\textwidth}X X@{}}
\toprule
\textbf{ID} & \textbf{Research question} & \textbf{Hypothesis, if applicable} \\
\midrule
RQ-01 & \placeholder{Focused and answerable research question} & \placeholder{H1 or Not applicable} \\
RQ-02 & \placeholder{Additional research question, if required} & \placeholder{H2 or Not applicable} \\
\bottomrule
\end{tabularx}
\end{table}

\section{Objectives}
\begin{table}[H]
\centering
\caption{Project objectives}
\label{tab:objectives}
\begin{tabularx}{\textwidth}{@{}p{0.14\textwidth}X@{}}
\toprule
\textbf{ID} & \textbf{Objective} \\
\midrule
OBJ-01 & Enable Customers and Contractors to create, discover, negotiate, reserve, and track lot-based surplus-material listings through a single marketplace. \\
OBJ-02 & Classify the four scoped material categories, detect selected visible defects, assess condition, check description--image consistency, and estimate a fair surplus price. \\
OBJ-03 & Provide human-in-the-loop review for uncertain listings and a contextual assistant that answers buyers using only available listing information. \\
\bottomrule
\end{tabularx}
\end{table}

\section{Proposed Approach}
The project will develop a marketplace with Admin, Customer, and Contractor
roles. Its AI pipeline uses CLIP zero-shot classification, material-specific YOLO
defect detection, a condition engine, description--image consistency checks, and
XGBoost price estimation. A decision engine auto-publishes clean listings and
sends uncertain or Mixed-condition listings to an Admin review queue. The core
platform provides lot-based listings, location and category search, real-time
chat, negotiation, ratings, saved-search alerts, reservation, and pickup-status
tracking.

\section{Scope}
\subsection{In Scope}

\begin{table}[H]
\centering
\caption{In-scope items}
\label{tab:in-scope}
\begin{tabularx}{\textwidth}{@{}p{0.14\textwidth}X p{0.20\textwidth}@{}}
\toprule
\textbf{Scope ID} & \textbf{Included item} & \textbf{Objective(s)} \\
\midrule
SCP-01 & Lot-based listings with quantity, description, price, and 3--5 representative photographs. & OBJ-01 \\
SCP-02 & Customer/Contractor search, negotiation, reservation, and pickup-status tracking. & OBJ-01 \\
SCP-03 & AI classification, defect detection, condition grading, consistency checking, and price estimation for the four scoped categories. & OBJ-02 \\
SCP-04 & Admin review queue with approval, rejection, edit requests, threshold tuning, and AI override. & OBJ-03 \\
SCP-05 & Location, price, and category search; real-time chat and negotiation; ratings, saved-search alerts, reservation, and pickup tracking. & OBJ-01 \\
\bottomrule
\end{tabularx}
\end{table}

\subsection{Out of Scope \applicability{If applicable}}
\begin{table}[H]
\centering
\caption{Out-of-scope items}
\label{tab:out-of-scope}
\begin{tabularx}{\textwidth}{@{}p{0.14\textwidth}X X@{}}
\toprule
\textbf{ID} & \textbf{Excluded item} & \textbf{Reason} \\
\midrule
EXC-01 & Integration with the university's central timetable system. & The external interface is unavailable during the FYP period. \\
EXC-02 & Hidden or internal defects and engineering-grade structural inspection. & The project targets only visible defects supported by validated datasets. \\
EXC-03 & Real payment processing and contractor/business financial verification. & Payment compliance and financial verification are outside the project goals. \\
EXC-04 & Guaranteeing the condition of every item in a large lot. & AI assessment applies only to uploaded representative samples. \\
EXC-05 & Automated dispute resolution and delivery. & Buyers and sellers arrange pickup; disputes remain manual Admin work. \\
\bottomrule
\end{tabularx}
\end{table}

\section{Users and Stakeholders}
\begin{table}[H]
\centering
\caption{Users and stakeholders}
\label{tab:stakeholders}
\begin{tabularx}{\textwidth}{@{}p{0.22\textwidth}p{0.18\textwidth}X X@{}}
\toprule
\textbf{Stakeholder} & \textbf{Role} & \textbf{Need or interest} & \textbf{Design influence} \\
\midrule
Contractor/Seller & Primary user & Create lots, negotiate, and sell surplus material & Lot-based listing and seller tools \\
Customer/Buyer & Primary user & Find trustworthy material and reserve lots & Location, price, category search and AI assistant \\
Admin & Moderator & Review flagged listings, manage users and categories, and override AI decisions & Review queue, thresholds, dashboard, and moderation controls \\
\bottomrule
\end{tabularx}
\end{table}

% =============================================================================

% This file is included by uet_thesis.tex.
% Replace blue placeholders and remove example text before submission.

\chapter{Requirements Specification}
\label{chap:requirements}
% =============================================================================

\section{Assumptions, Constraints, and Dependencies}
\begin{table}[H]
\centering
\caption{Assumptions, constraints, and dependencies}
\label{tab:acd}
\begin{tabularx}{\textwidth}{@{}p{0.13\textwidth}p{0.18\textwidth}X X@{}}
\toprule
\textbf{ID} & \textbf{Type} & \textbf{Statement} & \textbf{Implication} \\
\midrule
ACD-01 & Assumption & Users have access to a modern web browser. & A separate desktop application is not required. \\
ACD-02 & Constraint & AI assessment covers only four material categories and validated visible defect classes. & The system must disclose its limits and avoid claiming exhaustive inspection. \\
ACD-03 & Dependency & Public defect datasets, published new-material prices, and a small interview/validation sample. & Model coverage and price depreciation estimates depend on data availability and quality. \\
\bottomrule
\end{tabularx}
\end{table}

\section{Functional Requirements \applicability{Development, Hybrid, Game; Research Software Prototype}}
\begin{table}[H]
\centering
\caption{Functional requirements}
\label{tab:functional-requirements}
\begin{tabularx}{\textwidth}{@{}p{0.12\textwidth}X p{0.15\textwidth}p{0.20\textwidth}@{}}
\toprule
\textbf{ID} & \textbf{Requirement} & \textbf{Priority} & \textbf{Source} \\
\midrule
FR-01 & The system shall allow a seller to create a lot listing with quantity, description, price, and 3--5 representative photographs. & Must & Contractor/Seller \\
FR-02 & The system shall classify submitted listing photos into Bricks, Steel/Rebar, Tiles/Marble, or Wood. & Must & AI pipeline \\
FR-03 & The system shall detect configured visible defects and derive a Good, Minor, Major, Mixed, or Not Assessed condition grade. & Must & AI pipeline \\
FR-04 & The system shall route uncertain, inconsistent, or Mixed-condition listings to the Admin review queue. & Must & Admin \\
FR-05 & The system shall allow Admins to approve, reject, request edits, or override an AI listing decision. & Must & Admin \\
FR-06 & The system shall allow users to search listings by location, price, and category. & Must & Customer/Contractor \\
FR-07 & The system shall support real-time chat, negotiation, and lot reservation. & Must & Customer/Contractor \\
FR-08 & The system shall answer buyer questions using only the listing's structured data and return an explicit unavailable response when information is missing. & Should & Customer/Buyer \\
\bottomrule
\end{tabularx}
\end{table}

\section{Non-Functional Requirements \applicability{Development, Hybrid, Game; Research Software Prototype}}
\begin{table}[H]
\centering
\caption{Non-functional requirements}
\label{tab:nonfunctional-requirements}
\begin{tabularx}{\textwidth}{@{}p{0.12\textwidth}p{0.18\textwidth}X p{0.15\textwidth}@{}}
\toprule
\textbf{ID} & \textbf{Attribute} & \textbf{Measurable requirement} & \textbf{Priority} \\
\midrule
NFR-01 & Reliability & AI components shall be decoupled so that a temporarily unavailable AI service does not prevent a submitted listing from remaining valid and purchasable. & Must \\
NFR-02 & Transparency & The system shall disclose that AI assessment applies to representative samples and does not guarantee every unphotographed item. & Must \\
\bottomrule
\end{tabularx}
\end{table}

\section{Data and External Interface Requirements \applicability{If applicable}}
\begin{table}[H]
\centering
\caption{Data and external interface requirements}
\label{tab:data-interface-requirements}
\begin{tabularx}{\textwidth}{@{}p{0.12\textwidth}p{0.19\textwidth}X X@{}}
\toprule
\textbf{ID} & \textbf{Type} & \textbf{Requirement} & \textbf{Source or system} \\
\midrule
DIR-01 & Data & The system shall store each lot with quantity, material, photographs, condition, detected defects, price, location, and status. & Listing process \\
EIR-01 & External API & The core API shall communicate asynchronously with the Python/FastAPI AI service for classification, defect detection, consistency, and price results. & AI service \\
EIR-02 & Real-time service & The system shall exchange chat and negotiation messages through the SignalR service. & SignalR \\
EIR-03 & Object storage & The system shall store listing photographs in AWS S3 and associate them with their lot listing. & AWS S3 \\
\bottomrule
\end{tabularx}
\end{table}

\section{Legal, Ethical, Privacy, and Safety Requirements \applicability{If applicable}}
\begin{table}[H]
\centering
\caption{Legal, ethical, privacy, and safety requirements}
\label{tab:leps-requirements}
\begin{tabularx}{\textwidth}{@{}p{0.12\textwidth}p{0.18\textwidth}X X@{}}
\toprule
\textbf{ID} & \textbf{Category} & \textbf{Requirement or obligation} & \textbf{Applicable authority or source} \\
\midrule
PR-01 & Privacy & Personal data shall be collected only for stated project purposes and accessed only by authorized roles. & Project privacy requirement \\
ER-01 & Ethics & Interview and questionnaire participants shall provide informed consent, and collected responses shall be used only for requirements and evaluation. & Supervisor/institutional guidance \\
\bottomrule
\end{tabularx}
\end{table}

% =============================================================================

% This file is included by uet_thesis.tex.
% Replace blue placeholders and remove example text before submission.

\chapter{Literature Review and Related Work}
\label{chap:literature}
% =============================================================================

\section{Review Method}
\begin{guidance}
Briefly state how related work was identified and selected. Report the sources
searched, key search terms, publication period, and principal inclusion and
exclusion criteria. Follow departmental requirements for the minimum number and
type of sources.
\end{guidance}

\begin{table}[H]
\centering
\caption{Review method summary}
\label{tab:review-method}
\begin{tabularx}{\textwidth}{@{}p{0.27\textwidth}X@{}}
\toprule
\textbf{Item} & \textbf{Description} \\
\midrule
Sources searched & \placeholder{e.g., IEEE Xplore, ACM Digital Library, Scopus, official product sources} \\
Search terms & \placeholder{Main keywords and combinations} \\
Coverage period & \placeholder{Years covered and search date} \\
Selection criteria & \placeholder{Principal inclusion and exclusion criteria} \\
\bottomrule
\end{tabularx}
\end{table}

\section{Related Systems and Solutions \applicability{Development, Hybrid, Game}}
\begin{table}[H]
\centering
\caption{Related systems and solutions}
\label{tab:related-systems}
\begin{tabularx}{\textwidth}{@{}p{0.20\textwidth}X X X@{}}
\toprule
\textbf{System and citation} & \textbf{Relevant features} & \textbf{Strengths} & \textbf{Limitations} \\
\midrule
OLX Pakistan & General classifieds and communication & Familiar marketplace workflow and chat & No automated condition, price, or description verification \\
Facebook Marketplace & General listings and discovery & Large user base and simple listing flow & No structured construction-material categories or photo-based assessment \\
OfferUp & Consumer resale marketplace & Supports peer-to-peer resale & Not specialized for construction material and lacks domain verification \\
Zarea.pk & Construction-material marketplace & Focuses on construction-material discovery & Primarily new materials from dealers rather than verified peer-to-peer surplus \\
\bottomrule
\end{tabularx}
\end{table}

\section{Literature Synthesis \applicability{Research, Hybrid; other projects where applicable}}
\begin{guidance}
Organize scholarly literature by themes, methods, findings, or disagreements.
Compare and connect studies instead of presenting an isolated summary of each
paper. Identify limitations that directly inform the project.
\end{guidance}

The proposed technical foundation combines three established themes: CLIP
zero-shot image--text matching for material classification, YOLO-based object
detection for visible defects, and XGBoost gradient-boosted regression for
structured price estimation. Public defect datasets support model development,
but their inspection imagery may differ from marketplace phone photographs.
Therefore, the project uses a separate real-photo validation set and limits its
claims to defect classes that are present and validated in the selected datasets.

\section{Comparative Analysis}
\begin{table}[H]
\centering
\caption{Comparative analysis}
\label{tab:comparative-analysis}
\begin{tabularx}{\textwidth}{@{}p{0.26\textwidth}X X X@{}}
\toprule
\textbf{System or study} & \textbf{Criterion 1} & \textbf{Criterion 2} & \textbf{Criterion 3} \\
\midrule
OLX Pakistan & Manual moderation & General listings & No condition or price verification \\
Facebook Marketplace & User-reported information & General categories & No domain-specific defect assessment \\
\textbf{Proposed project} & \textbf{AI-gated plus Admin review} & \textbf{Four construction-material categories} & \textbf{Condition, description, and price checks} \\
\bottomrule
\end{tabularx}
\end{table}

\section{Gap Analysis and Proposed Contribution}
Existing marketplaces provide discovery and communication but do not verify
construction-material condition, description accuracy, or surplus-price fairness
before publication. surplusSpace contributes a bounded, human-in-the-loop
marketplace that combines material classification, visible-defect detection,
condition grading, description--image consistency, and price estimation, while
retaining Admin authority for uncertain cases.

% =============================================================================

% This file is included by uet_thesis.tex.
% Replace blue placeholders and remove example text before submission.

\chapter{Design Specification}
\label{chap:design}
% =============================================================================

\section{System Context, Architecture, and Interfaces \applicability{Development, Hybrid, Game; Research Software Prototype}}
The system boundary contains the Flutter web/Android client, the ASP.NET Core
marketplace API, the PostgreSQL/PostGIS database, Redis, SignalR, and the
asynchronous Python/FastAPI AI service. Customers and Contractors submit and
discover lots, communicate, negotiate, and reserve them. Admins manage users,
categories, thresholds, flagged listings, and AI overrides. AWS S3 stores listing
photographs. The AI service returns classification, defect, condition,
consistency, and price signals to the decision engine.

\begin{table}[H]
\centering
\caption{Major interfaces}
\label{tab:major-interfaces}
\begin{tabularx}{\textwidth}{@{}p{0.16\textwidth}p{0.22\textwidth}X X@{}}
\toprule
\textbf{ID} & \textbf{Interface} & \textbf{Purpose and exchanged information} & \textbf{Protocol or format} \\
\midrule
INT-01 & Flutter client--ASP.NET Core API & Authentication, listings, search, reservations, and Admin operations & REST/JSON \\
INT-02 & ASP.NET Core--FastAPI AI service & Listing photos and metadata exchanged for AI assessment results & Asynchronous API \\
INT-03 & ASP.NET Core--SignalR & Real-time chat and negotiation messages & SignalR \\
INT-04 & Platform--PostgreSQL/PostGIS and Redis & Persistent marketplace data, geospatial queries, cache, and presence & Database/cache protocols \\
INT-05 & Platform--AWS S3 & Listing-photo upload and retrieval & Object storage API \\
\bottomrule
\end{tabularx}
\end{table}

\section{Research Methodology and Experimental Design \applicability{Research, Hybrid}}
\subsection{Study or Experimental Design}
The project will evaluate the integrated system using held-out image data and
Admin review decisions. Each submitted lot passes through the same classification,
defect, condition, consistency, and price stages. Results are compared with
validated labels, sourced reference prices, and human review outcomes.

\subsection{Population, Sample, Participants, or Dataset}
The AI pipeline will use existing public building-defect datasets, including
YOLO-format datasets, SDNET2018, and MBDD2025, supplemented by a small
team-collected validation set of marketplace-style phone photographs. New-material
prices will be collected from published Pakistani construction-pricing sources.
Surplus depreciation will be informed by structured interviews with 5--10 Lahore
dealers and contractors. User requirements will be validated through interviews
with 4--6 contractors/suppliers and a brief questionnaire.

\subsection{Variables and Data Collection Techniques}

\begin{table}[H]
\centering
\caption{Variables and data collection}
\label{tab:variables}
\begin{tabularx}{\textwidth}{@{}p{0.16\textwidth}p{0.23\textwidth}X X@{}}
\toprule
\textbf{ID} & \textbf{Variable or construct} & \textbf{Operational definition or measure} & \textbf{Collection source or instrument} \\
\midrule
VAR-01 & \placeholder{Variable or construct} & \placeholder{How it is defined or measured} & \placeholder{Instrument, sensor, dataset field, interview, etc.} \\
VAR-02 & Material classification accuracy & Correct predicted category across Bricks, Steel/Rebar, Tiles/Marble, and Wood & Held-out image set \\
VAR-03 & Defect detection performance & Precision and recall for each configured defect class & Labelled public datasets and real-photo validation set \\
VAR-04 & Price-estimate error & MAE and RMSE against reference prices & Published Pakistani price sources \\
VAR-05 & Decision agreement & Agreement between automated publish/flag result and Admin decision & Admin review outcomes \\
\bottomrule
\end{tabularx}
\end{table}

\subsection{Analysis and Evaluation Protocol}
Evaluation will report material-classification accuracy and confusion patterns,
defect-detection precision and recall per defect class, agreement between the
automated publish/flag decision and Admin review decisions, and price-estimate
MAE/RMSE against sourced reference prices. The real-photo validation set will be
kept separate from training data.



\section{Behavioural and Interaction Design \applicability{Development, Hybrid, Game; Research if applicable}}
\begin{guidance}
Include only the models needed to explain the project's behaviour and interaction.
Each UML diagram must use standard UML notation~\citep{omguml251}, remain readable, and be followed by a
brief explanation of important flows, alternatives, or design decisions.
\end{guidance}

\subsection{Use-Case Model}
The principal use cases are Create Lot Listing, Search Listings, Chat and
Negotiate, Reserve Lot, Review Flagged Listing, Override AI Decision, and Ask the
AI Listing Assistant.

\begin{table}[H]
\centering
\caption{Use-case specification}
\label{tab:use-case-specification}
\begin{tabularx}{\textwidth}{@{}p{0.24\textwidth}X@{}}
\toprule
\textbf{Field} & \textbf{Description} \\
\midrule
Use-case ID and name & UC-01: Create and verify lot listing \\
Primary actor & Contractor/Seller \\
Preconditions & Authenticated account and lot details with 3--5 photographs \\
Trigger & Seller submits a listing \\
Main flow & Store lot; run AI pipeline; publish if checks pass or queue for Admin review \\
Alternative or exception flows & Low confidence, disagreement, or Mixed condition routes to Admin \\
Postconditions & Listing is published, rejected, or awaiting review \\
Related requirements & FR-01--FR-05, NFR-01--NFR-02 \\
\bottomrule
\end{tabularx}
\end{table}

\subsection{Activity Model}
\begin{guidance}
Model an important workflow using actions, control flows, decisions with guards,
parallel paths where required, and clear initial and final nodes. Use swimlanes when
responsibility changes are important.
\end{guidance}
\diagramplaceholder{Activity diagram}{Example workflow: submit request, check availability, approve or reject, and notify the student.}

\subsection{State-Machine Model \applicability{State-dependent systems and games}}
\begin{guidance}
Model the lifecycle of one important stateful entity. Show states, triggering
events, guards where needed, actions, and initial and final states.
\end{guidance}
\diagramplaceholder{State-machine diagram}{Example entity: Appointment; states may include Requested, Approved, Rescheduled, Completed, and Cancelled.}

\subsection{Sequence Model}
\begin{guidance}
Show how an important scenario is realized through time-ordered messages among
actors, system components, and external services. Use activation bars and combined
fragments such as \texttt{alt}, \texttt{opt}, or \texttt{loop} where appropriate.
\end{guidance}
\diagramplaceholder{Sequence diagram}{Example scenario: a student submits an appointment request and receives confirmation.}

\section{Structural Design \applicability{Development, Hybrid, Game; Research if applicable}}
\begin{guidance}
Describe the static structure needed to realize the requirements. Keep models at a
level appropriate to the project and avoid including generated framework classes
that do not communicate a design decision.
\end{guidance}

\subsection{Class or Domain Model}
\begin{guidance}
Show important domain concepts or software classes, relevant attributes and
operations, relationships, multiplicities, and inheritance where justified.
\end{guidance}
\diagramplaceholder{Class or domain diagram}{Example classes: Student, Adviser, Appointment, AvailabilitySlot, and Notification.}

\subsection{Component Model}
\begin{guidance}
Show deployable or replaceable software components, their provided and required
interfaces, and dependencies.
\end{guidance}
\diagramplaceholder{Component diagram}{Example components: Web Client, Appointment Service, Authentication Service, and Data Store.}

\subsection{Package Model \applicability{If required}}
\begin{guidance}
Use a package diagram when the organization and dependencies of modules or layers
need clarification. Omit it when the component or class model already communicates
the structure adequately.
\end{guidance}
\diagramplaceholder{Package diagram}{Example packages: presentation, application, domain, and infrastructure.}

\subsection{Deployment Model \applicability{Deployed, distributed, networked, or multi-device systems}}
\begin{guidance}
Show execution nodes, devices, environments, deployed artifacts, and communication
paths. Distinguish physical or virtual nodes from software components.
\end{guidance}
\diagramplaceholder{Deployment diagram}{Example nodes: user device, application server, database server, and external notification service.}

\section{Data Design and Management}
\label{sec:data-design}

\subsection{Database, Schema, and Persistence Design \applicability{Development, Hybrid, Game, if applicable}}
\begin{guidance}
Describe persistent entities, identifiers, relationships, integrity constraints,
and important indexing or retention decisions. Use an ER diagram, schema, or other
notation appropriate to the selected database model.
\end{guidance}
\diagramplaceholder{Data model or schema}{Example entities: User, AvailabilitySlot, Appointment, and NotificationLog.}

\subsection{Research Data Management \applicability{Research, Hybrid}}
\begin{guidance}
Explain how research data will be collected or acquired, named, stored, protected,
anonymized or pseudonymized, versioned, retained, shared, and disposed of. Identify
access controls and responsible team members.
\end{guidance}

\begin{table}[H]
\centering
\caption{Research data management plan}
\label{tab:data-management}
\begin{tabularx}{\textwidth}{@{}p{0.22\textwidth}X X@{}}
\toprule
\textbf{Data category} & \textbf{Management method} & \textbf{Protection, retention, or access rule} \\
\midrule
\placeholder{Raw/processed data, code, consent records, model artifacts} & \placeholder{Format, location, naming, and versioning} & \placeholder{Access, retention, sharing, and disposal} \\
\bottomrule
\end{tabularx}
\end{table}

\section{Algorithm or Model Design \applicability{Research, Hybrid; Development or Game if a core algorithm or AI model is used}}
\begin{table}[H]
\centering
\caption{Core algorithm or model}
\label{tab:algorithm-model}
\begin{tabularx}{\textwidth}{@{}p{0.18\textwidth}X@{}}
\toprule
\textbf{Item} & \textbf{Description} \\
\midrule
Name and purpose & CLIP, YOLO, condition engine, consistency checker, and XGBoost price model for listing verification. \\
Inputs and outputs & Representative listing photographs, description, quantity, location, and reference price; material, defects, condition, consistency result, and estimated surplus price. \\
Processing steps & Classify material; detect visible defects; aggregate photo results into a lot grade; compare description and image signals; estimate price; publish or flag using decision thresholds. \\
Assumptions and limits & Assessment covers representative photos and only defect classes present and validated in available datasets. \\
Related requirements & FR-02--FR-04, NFR-02 \\
\bottomrule
\end{tabularx}
\end{table}

\section{User Interface and Accessibility Design \applicability{User-facing Development, Hybrid, and Game projects; Research prototypes if applicable}}
\begin{guidance}
Present the navigation structure and key wireframes or prototypes. Explain how the
design supports usability, accessibility, responsive layout, error prevention,
feedback, and consistency. Do not include every minor screen.
\end{guidance}
\diagramplaceholder{Navigation structure or user flow}{Example: Sign In, Dashboard, Available Slots, Appointment Details, and Notifications.}
\diagramplaceholder{Key wireframe or prototype}{Annotate important controls, content, validation, feedback, and accessibility decisions.}

\section{Security, Privacy, and Safety Design \applicability{If applicable}}
\begin{guidance}
Explain the design controls used to satisfy the corresponding requirements in
Chapter~\ref{chap:requirements}. Address relevant threats, authentication,
authorization, data protection, logging, secure failure, privacy, and safety. Do
not make general claims such as ``the system is secure'' without specific controls.
\end{guidance}

\begin{table}[H]
\centering
\caption{Security, privacy, and safety controls}
\label{tab:security-controls}
\begin{tabularx}{\textwidth}{@{}p{0.16\textwidth}X X p{0.19\textwidth}@{}}
\toprule
\textbf{Design ID} & \textbf{Risk or requirement} & \textbf{Design control} & \textbf{Related requirement(s)} \\
\midrule
SEC-01 & Unauthorized access to appointment records & Role-based authorization is checked at the service layer before records are returned or modified. & FR-XX, NFR-XX \\
SEC-02 & \placeholder{Threat, privacy concern, or safety hazard} & \placeholder{Specific design control} & \placeholder{Requirement IDs} \\
\bottomrule
\end{tabularx}
\end{table}

\section{Design Decisions and Rationale}
\begin{guidance}
Record significant choices, alternatives considered, and reasons for selection.
Focus on decisions that materially affect quality, cost, risk, maintainability, or
project feasibility. Do not repeat the tools list from Section~\ref{sec:tools}.
\end{guidance}

\begin{table}[H]
\centering
\caption{Design decisions}
\label{tab:design-decisions}
\begin{tabularx}{\textwidth}{@{}p{0.14\textwidth}X X X@{}}
\toprule
\textbf{ID} & \textbf{Decision} & \textbf{Alternatives considered} & \textbf{Rationale and trade-off} \\
\midrule
DD-01 & \placeholder{Selected design decision} & \placeholder{Feasible alternatives} & \placeholder{Evidence-based reason and accepted limitation} \\
\bottomrule
\end{tabularx}
\end{table}

% =============================================================================

% This file is included by uet_thesis.tex.
% Replace blue placeholders and remove example text before submission.

\chapter{Planning, Risk, and Traceability}
\label{chap:planning}
% =============================================================================

\section{Design Traceability Matrix}
{\footnotesize
\begin{longtable}{@{}p{0.12\textwidth}p{0.11\textwidth}p{0.16\textwidth}p{0.23\textwidth}p{0.24\textwidth}@{}}
\caption{Design traceability matrix}\label{tab:traceability}\\
\toprule
\textbf{Objective} & \textbf{Scope} & \textbf{Requirement} & \textbf{Design element} & \textbf{Testing-document reference} \\
\midrule
\endfirsthead
\toprule
\textbf{Objective} & \textbf{Scope} & \textbf{Requirement} & \textbf{Design element} & \textbf{Testing-document reference} \\
\midrule
\endhead
OBJ-01 & SCP-01 & FR-01 & UC-01; Lot Listing; Listing model & Testing document: listing creation \\
OBJ-01 & SCP-05 & FR-06--FR-07 & UC-02; search and chat components; reservation model & Testing document: discovery and reservation \\
OBJ-02 & SCP-03 & FR-02--FR-03 & AI pipeline; condition engine; AI model design & Testing document: classification and defect evaluation \\
OBJ-03 & SCP-04 & FR-04--FR-05 & Admin review queue; decision engine & Testing document: decision agreement and override \\
\bottomrule
\end{longtable}
}

\section{Verification and Validation Strategy}
\begin{table}[H]
\centering
\caption{Verification and validation strategy}
\label{tab:vv-strategy}
\begin{tabularx}{\textwidth}{@{}p{0.20\textwidth}X X p{0.20\textwidth}@{}}
\toprule
\textbf{Target} & \textbf{Method} & \textbf{Evidence} & \textbf{Responsible role} \\
\midrule
Functional requirements & Functional and integration testing & Approved test results in the testing document & Whole team \\
Stakeholder needs & Demonstration and structured interviews/questionnaire & Feedback and validation record & Whole team and Admin \\
Research findings & Accuracy, precision/recall, decision agreement, MAE/RMSE evaluation & Model and evaluation outputs & AI/ML lead \\
\bottomrule
\end{tabularx}
\end{table}

\section{Testing-Document Reference}
\begin{guidance}
Identify the separate testing document that contains the detailed acceptance
criteria, test cases, datasets, procedures, results, and defect records. Keep its
version consistent with this design document.
\end{guidance}

\begin{table}[H]
\centering
\caption{Software testing document reference}
\label{tab:testing-reference}
\begin{tabularx}{\textwidth}{@{}p{0.25\textwidth}X@{}}
\toprule
\textbf{Item} & \textbf{Reference} \\
\midrule
Document title & \placeholder{Project Software Testing Document} \\
Version and date & \placeholder{Version; Day Month Year} \\
Repository or submission location & \placeholder{Controlled path, repository location, or submission reference} \\
Responsible team member(s) & \placeholder{Name(s) and role(s)} \\
\bottomrule
\end{tabularx}
\end{table}

\section{Risk Assessment}
\begin{table}[H]
\centering
\caption{Risk register}
\label{tab:risk-register}
\small
\begin{tabularx}{\textwidth}{@{}p{0.11\textwidth}X p{0.14\textwidth}p{0.11\textwidth}X p{0.14\textwidth}@{}}
\toprule
\textbf{ID} & \textbf{Risk} & \textbf{Likelihood} & \textbf{Impact} & \textbf{Response} & \textbf{Owner} \\
\midrule
RSK-01 & AI service becomes unavailable. & Medium & High & Keep AI services decoupled so a listing remains valid and purchasable without a price suggestion. & Backend team \\
RSK-02 & Public defect datasets differ from marketplace phone photos. & Medium & High & Validate on a small team-collected real-photo set before relying on model results. & AI/ML lead \\
RSK-03 & Surplus-price depreciation lacks published data. & Medium & Medium & Ground assumptions in interviews with 5--10 Lahore dealers and contractors and refine them during evaluation. & AI/ML lead \\
\bottomrule
\end{tabularx}
\end{table}

\section{Project Management Approach and Schedule}
\begin{table}[H]
\centering
\caption{Milestones and schedule}
\label{tab:schedule}
\begin{tabularx}{\textwidth}{@{}p{0.14\textwidth}X p{0.20\textwidth}p{0.19\textwidth}p{0.16\textwidth}@{}}
\toprule
\textbf{ID} & \textbf{Milestone or deliverable} & \textbf{Target date} & \textbf{Dependency} & \textbf{Owner} \\
\midrule
MS-01 & Core platform and listing infrastructure & Month 1--4 & None & Backend/frontend teams \\
MS-02 & AI classification and defect pipeline & Month 3--7 & MS-01 & AI/ML team \\
MS-03 & Admin, chat, chatbot, and full integration & Month 6--9 & MS-01; MS-02 & Whole team \\
MS-04 & Testing, evaluation, and final release & Month 10--12 & MS-03 & Whole team \\
\bottomrule
\end{tabularx}
\end{table}

\section{Team Roles and Responsibilities}
\begin{table}[H]
\centering
\caption{Team roles and responsibilities}
\label{tab:team-roles}
\begin{tabularx}{\textwidth}{@{}p{0.24\textwidth}p{0.24\textwidth}X@{}}
\toprule
\textbf{Member} & \textbf{Role} & \textbf{Principal responsibilities} \\
\midrule
Hafsa Shahid (2023-CS-708) & Frontend developer & Flutter listing, search, chat, and Admin dashboard interfaces \\
Fatima Irfan (2023-CS-715) & Backend developer & REST API, PostgreSQL/PostGIS, RBAC, listings, search, reservations, and trust layer \\
Daniyal Jumshaid (2023-CS-734) & Real-time and AI integration developer & SignalR chat, FastAPI integration, decision engine, and Admin review backend \\
Noor Fatima Malik (2023-CS-739) & AI/ML lead & YOLO, CLIP, XGBoost, chatbot context, and system evaluation \\
\bottomrule
\end{tabularx}
\end{table}

\section{Source-Code, Data, and Artifact Management}
\begin{guidance}
State how code, documents, datasets, models, diagrams, builds, and other artifacts
will be named, versioned, reviewed, backed up, licensed, and accessed. Identify how
sensitive data and credentials will be kept out of the source repository.
\end{guidance}

\begin{table}[H]
\centering
\caption{Artifact management}
\label{tab:artifact-management}
\begin{tabularx}{\textwidth}{@{}p{0.21\textwidth}X X X@{}}
\toprule
\textbf{Artifact} & \textbf{Storage and versioning} & \textbf{Access or protection} & \textbf{Backup or retention} \\
\midrule
Source code & \placeholder{Repository and branching/review approach} & \placeholder{Access control and secret handling} & \placeholder{Backup or release policy} \\
\placeholder{Data, models, and documents} & \placeholder{Location and versioning} & \placeholder{Access rule} & \placeholder{Retention rule} \\
\bottomrule
\end{tabularx}
\end{table}

\section{Tools and Technologies}
\label{sec:tools}
\begin{table}[H]
\centering
\caption{Tools and technologies}
\label{tab:tools-technologies}
\begin{tabularx}{\textwidth}{@{}p{0.24\textwidth}p{0.20\textwidth}X X@{}}
\toprule
\textbf{Tool or technology} & \textbf{Version} & \textbf{Purpose} & \textbf{Project use} \\
\midrule
Flutter (Dart) & 3.x & Cross-platform UI & Web and Android marketplace client \\
ASP.NET Core (C\#) & 8 & Backend API framework & Core marketplace services \\
PostgreSQL/PostGIS & Current project version & Relational and geospatial database & Listings, users, locations, and search \\
SignalR and Redis & Current project versions & Real-time communication and presence/cache & Chat, negotiation, and presence \\
Python/FastAPI & Current project versions & AI service API & Asynchronous AI pipeline integration \\
PyTorch, OpenCV, CLIP, YOLO & Current project versions & Computer vision and classification & Material classification and defect detection \\
XGBoost and scikit-learn & Current project versions & Tabular regression and evaluation & Surplus-price estimation \\
AWS S3 & Current service version & Object storage & Listing photographs \\
Swagger, Postman, k6/JMeter & Current versions & API and load testing & Functional, API, concurrency, and load testing \\
Git and GitHub & Current versions & Version control & Source and artifact management \\
\bottomrule
\end{tabularx}
\end{table}


\addtocontents{toc}{\vspace{2em}}
\appendix
% This file is included by uet_thesis.tex.
% Replace blue placeholders and remove example text before submission.

\chapter{Definitions, Acronyms, and Domain Terms}
\label{app:terms}
\begin{guidance}
Define abbreviations and specialized terms that a reader may not know. Do not list
common computing terms unless they have a project-specific meaning.
\end{guidance}

\begin{longtable}{@{}p{0.27\textwidth}p{0.68\textwidth}@{}}
\caption{Definitions, acronyms, and domain terms}\label{tab:terms}\\
\toprule
\textbf{Term or acronym} & \textbf{Definition} \\
\midrule
\endfirsthead
\toprule
\textbf{Term or acronym} & \textbf{Definition} \\
\midrule
\endhead
\placeholder{Term} & \placeholder{Clear project-specific definition} \\
\placeholder{Acronym} & \placeholder{Expanded form and meaning} \\
\bottomrule
\end{longtable}

% This file is included by uet_thesis.tex.
% Replace blue placeholders and remove example text before submission.

\chapter{Supporting Specifications \applicability{If applicable}}
\begin{guidance}
Place lengthy supporting material here only when it is necessary for understanding
or reproducing the design, such as API contracts, data dictionaries, detailed
pseudocode, questionnaires, or supplementary diagrams. Refer to every included
item from the main document.
\end{guidance}

\placeholder{Insert supporting specification or remove this appendix.}


\addtocontents{toc}{\vspace{2em}}
\backmatter
\label{References}
\bibliographystyle{plainnat}
\bibliography{references}

\end{document}
